% Part: first-order-logic
% Chapter: tableaux
% Section: provability-consistency

\documentclass[../../../include/open-logic-section]{subfiles}

\begin{document}

\iftag{FOL}
      {\olfileid[pt]{fol}{tab}{prv}}
      {\olfileid[pt]{pl}{tab}{prv}}

\olsection{\usetoken{S}{derivability} e Consistência}

Agora estabeleceremos várias propriedades da relação de
!!{derivability}. Elas são interessantes por si mesmas, mas cada uma
desempenhará um papel na demonstração do teorema da completude.

\begin{prop}\ollabel{prop:provability-contr}
  Se $\Gamma \Proves !A$ e $\Gamma \cup \{!A\}$ é inconsistente,
  então $\Gamma$ é inconsistente.
\end{prop}

\begin{proof}
Existem conjuntos finitos $\Gamma_0 = \{!B_1, \dots, !B_n\}$ e $\Gamma_1
  =\{!C_1, \dots, !C_n\} \subseteq \Gamma$ tais que
  \begin{align*}
    \{\sFmla{\False}{!A}, &
    \sFmla{\True}{!B_1}, \dots, \sFmla{\True}{!B_n}\} \\
    \{\sFmla{\True}{!A}, &
    \sFmla{\True}{!C_1}, \dots, \sFmla{\True}{!C_m}\}
  \end{align*}
  têm !!{tableau}s fechados. Usando a regra \Cut{} em $!A$, podemos
  combiná-los em um único !!{tableau} fechado que mostra que $\Gamma_0
  \cup \Gamma_1$ é inconsistente. Como $\Gamma_0
  \subseteq \Gamma$ e $\Gamma_1 \subseteq \Gamma$, $\Gamma_0 \cup
  \Gamma_1 \subseteq \Gamma$, logo $\Gamma$~é inconsistente.
\end{proof}

\begin{prop}
\ollabel{prop:prov-incons}
$\Gamma \Proves !A$ se, e somente se, $\Gamma \cup \{\lnot !A\}$ é
inconsistente.
\end{prop}

\begin{proof}
Primeiro, suponha $\Gamma \Proves !A$, isto é, existe um
!!{tableau} fechado para
\[
\{\sFmla{\False}{!A},
\sFmla{\True}{!B_1}, \dots, \sFmla{\True}{!B_n}\}
\]
Usando a regra $\TRule{\True}{\lnot}$, ele pode ser transformado em um
!!{tableau} fechado para
\[
\{\sFmla{\True}{\lnot !A},
\sFmla{\True}{!B_1}, \dots, \sFmla{\True}{!B_n}\}.
\]

Por outro lado, se existe um !!{tableau} fechado para a segunda expressão,
podemos transformá-lo em um !!{tableau} fechado para a primeira removendo
toda fórmula que resulte da aplicação de \TRule{\True}{\lnot} à primeira
hipótese~$\sFmla{\True}{\lnot !A}$, bem como essa própria hipótese, e
acrescentando a hipótese $\sFmla{\False}{!A}$. De fato, se antes um ramo
estava fechado porque continha a conclusão de \TRule{\True}{\lnot}
aplicada a $\sFmla{\True}{\lnot !A}$, isto é, $\sFmla{\False}{!A}$, o
ramo correspondente no novo !!{tableau} também está fechado. Se um ramo
do tableau antigo estava fechado porque continha a hipótese
$\sFmla{\True}{\lnot !A}$, assim como $\sFmla{\False}{\lnot !A}$,
podemos transformá-lo em um ramo fechado aplicando
$\TRule{\False}{\lnot}$ a $\sFmla{\False}{\lnot !A}$ para obter
$\sFmla{\True}{!A}$. Isso fecha o ramo, pois acrescentamos
$\sFmla{\False}{!A}$ como hipótese.
\end{proof}

\begin{prob}
Demonstre que $\Gamma \Proves \lnot !A$ se, e somente se,
$\Gamma \cup \{!A\}$ é inconsistente.
\end{prob}

\begin{prop}\ollabel{prop:explicit-inc}
  Se $\Gamma \Proves !A$ e $\lnot !A \in \Gamma$, então $\Gamma$ é
  inconsistente.
\end{prop}

\begin{proof}
  Suponha $\Gamma \Proves !A$ e $\lnot !A \in \Gamma$. Então existem
  $!B_1$, \dots, $!B_n \in \Gamma$ tais que \
  \[
  \{\sFmla{\False}{!A}, \sFmla{\True}{!B_1}, \dots, \sFmla{\True}{!B_n}\}
  \]
  tem um tableau fechado. Substitua a hipótese \sFmla{\False}{!A} por
  \sFmla{\True}{\lnot !A} e insira, depois das hipóteses, a conclusão de
  \TRule{\True}{\lnot} aplicada a \sFmla{\False}{!A}. Qualquer
  !!{sentence} no !!{tableau} justificada por referência à linha~$1$ do
  !!{tableau} antigo passa a ser justificada por referência à linha~$n+1$.
  Portanto, se o !!{tableau} antigo era fechado, o novo também é. Ele mostra
  que $\Gamma$ é inconsistente, pois todas as hipóteses pertencem
  a~$\Gamma$.
\end{proof}

\begin{prop}\ollabel{prop:provability-exhaustive}
  Se $\Gamma \cup \{!A\}$ e $\Gamma \cup \{\lnot !A\}$ são ambos
  inconsistentes, então $\Gamma$ é inconsistente.
\end{prop}

\begin{proof}
  Se existem $!B_1$, \dots, $!B_n \in \Gamma$ e $!C_1$, \dots,
  $!C_m \in \Gamma$ tais que
  \begin{align*}
    \{\sFmla{\True}{!A}, &
    \sFmla{\True}{!B_1}, \dots, \sFmla{\True}{!B_n}\} \text{ e}\\
    \{\sFmla{\True}{\lnot !A}, &
    \sFmla{\True}{!C_1}, \dots, \sFmla{\True}{!C_m}\}
  \end{align*}
  têm !!{tableau}s fechados, podemos construir um único !!{tableau}
  combinado que mostra que $\Gamma$ é inconsistente, usando como hipóteses
  $\sFmla{\True}{!B_1}$, \dots, $\sFmla{\True}{!B_n}$ juntamente com
  $\sFmla{\True}{!C_1}$, \dots, $\sFmla{\True}{!C_m}$, seguidas de uma
  aplicação da regra \Cut{}. Isso produz dois ramos, um começando com
  $\sFmla{\True}{!A}$ e o outro com $\sFmla{\False}{!A}$.  
  
  No lado esquerdo, acrescente a parte do primeiro !!{tableau}
  abaixo de suas hipóteses. Aqui, toda aplicação de regra continua correta,
  pois cada uma das hipóteses do primeiro !!{tableau}, inclusive
  $\sFmla{\True}{!A}$, está disponível. Assim, todo ramo abaixo de
  $\sFmla{\True}{!A}$ se fecha. 
  
  No lado direito, acrescente a parte do segundo !!{tableau} abaixo de sua
  hipótese, removendo os resultados de quaisquer aplicações
  de~$\TRule{\True}{\lnot}$ a $\sFmla{\True}{\lnot !A}$. A conclusão de
  $\TRule{\True}{\lnot}$ aplicado a $\sFmla{\True}{\lnot !A}$ é
  $\sFmla{\False}{!A}$, que, não obstante, está disponível, pois é a
  conclusão da regra \Cut{} no lado direito do !!{tableau} combinado.

  Se um ramo do segundo tableau estava fechado porque continha a hipótese
  $\sFmla{\True}{\lnot !A}$ (que já não aparece como hipótese no
  !!{tableau} combinado), bem como $\sFmla{\False}{\lnot !A}$, podemos
  aplicar $\TRule{\False}{\lnot}$ a $\sFmla{\False}{\lnot !A}$ para obter
  $\sFmla{\True}{!A}$. Agora o ramo correspondente do !!{tableau} combinado
  também se fecha, pois contém a conclusão do lado direito da regra \Cut{},
  $\sFmla{\False}{!A}$. Se um ramo do segundo !!{tableau} se fechava por
  qualquer outro motivo, o ramo correspondente do !!{tableau} combinado
  também se fecha, pois quaisquer !!{signed formula}s diferentes de
  $\sFmla{\True}{\lnot !A}$ que ocorram no ramo do segundo !!{tableau}
  antigo também ocorrem no ramo correspondente do !!{tableau} combinado.
  \end{proof}

\end{document}
